% !TEX root = ../main.tex
\chapter{Analyse: run a chain over it}\label{ch:analyse}
Analyse is where you do the analysis itself. You take a \term{span}{span} of one
\term{channel}{channel}, pass it through a \term{chain}{chain} of \term{block}{blocks}, and look
at what every block did to it. This is the central skill of the app: every other workspace either
feeds Analyse (Explore, Library, Discovery) or reads what it produced (Review, Models).

\begin{plainterms}[A chain is a recipe]
Think of a kitchen recipe. The span is the raw ingredient. Each block is one step of the recipe:
wash, chop, simmer. Each step takes what the step before it handed over and hands something new to
the next. A step can only take an ingredient it knows how to handle: you cannot simmer a whole
uncut pumpkin in a teacup. In the app that rule is checked by the \emph{type} each block accepts
and emits (a signal, a set of spans, an encoding, and so on), and a chain that breaks it is marked
in red before you can run it.
\end{plainterms}

The workspace has three groups of pages.
\begin{itemize}
  \item \apppage{analyse.chain} and \apppage{analyse.block} are the working pair. The chain page
        shows every block as one row, each row drawing its result on one shared time axis; the
        block page opens one block, with its settings beside its output. These two pages are
        \pagestatus{live}: they run the real blocks on your recordings through the bridge.
  \item \apppage{analyse.glyphs} is a reference page listing the small picture each block carries.
  \item The Interrogation pages (\apppage{analyse.interrogation} and its two block pages) and the
        Training pages (\apppage{analyse.training} and its five block pages) are design previews of
        two special kinds of chain. Interrogation reads real families from the seed store but
        simulates its run; Training is example data throughout. Each section below says which
        parts are real.
\end{itemize}

The individual blocks are not explained here: each has a card in \cref{ch:blocks}, linked from
the text wherever a block is named.

% ---------------------------------------------------------------------------------------------
\pagesection{analyse.chain}{Chain}{Build a chain, or import a template}{live}
The chain page is where you assemble a chain, run it over a span and read every intermediate
result. The source (your span) is the top row; each block below it is one row, numbered
\texttt{01}, \texttt{02}, and so on; data flows from the top row down. Every row draws its own
result on the same time axis, labelled at the foot of the page (``all rows share this time
axis''), so a spike in row 01 sits directly above whatever row 03 made of it.

\pagefield{When you'd come here}
\begin{itemize}
  \item You sent a span from Explore (its \ui{Send span to Analyse} action) and want to analyse it.
  \item You want to try a saved \term{template}{template} on a new span, or tune one.
  \item You want to run one block on its own: that is a chain with one row.
  \item You want to start from saved \term{window-set}{window sets} rather than a span, to build
        the clustering chain that Models trains on.
\end{itemize}

\screenshot{pages/analyse/analyse.chain--live}{The page with the three-block example chain and no
source yet. Each row has a \ui{+ insert} pill above it; the footer chip names the chain's terminal
type and what it saves as.}

\pagefield{Doing things here}

\subsection*{Start a chain}
A fresh page opens on an empty canvas with a \ui{Start a chain} card. It offers:
\begin{enumerate}
  \item \ui{Import a template}: opens the templates list (see ``Import a template'' below).
  \item \ui{Add the first block}: opens the \ui{Insert a stage} window at position 01.
  \item \ui{Restore \ldots}: the last chain this browser had, if there was one. It is offered,
        never loaded unasked.
  \item \ui{Example \ldots}: a three-block example (baseline removal, then a matrix profile,
        then a threshold), useful for a first try.
\end{enumerate}

\subsection*{Choose the source}
The source chip in the toolbar names the span the chain runs over. Click it for the
\ui{Source} menu:
\begin{enumerate}
  \item \ui{Use the example span}: a two-hour reference span of \texttt{CH1\_A1} in
        \texttt{M2\_aug\_concat\_fs1.mat} (336--338\,h).
  \item \ui{Pick in Explore}: go to \apppage{explore.corpus} and send a span back.
  \item \ui{Start from a Window pool}: replace the span with saved window sets. The page swaps the
        chain for \emph{Window pool} \textrightarrow{} \emph{Trace shape} \textrightarrow{}
        \emph{Shape clustering} (\blockref{preprocessing.window_pool},
        \blockref{preprocessing.trace_shape}, \blockref{catalogue.shape_cluster}); your old chain
        is kept in the undo list. You tick the window sets on the pool block's own page.
\end{enumerate}
With no source, the source row says so and offers \ui{Go to Explore}, the example span and the
window pool. The held-out recording (\texttt{M4\_aug\_concat\_fs1.mat}) is never loaded: its
source row shows a lock card and \ui{Run chain} stays disabled.

\subsection*{Insert a stage}
\begin{enumerate}
  \item Click a \ui{+ insert} pill: the one above a row inserts there; \ui{+ insert · end of
        chain} appends.
  \item The \ui{Insert a stage} window opens. Its top line shows the chain as a ribbon with the
        new stage in place, and three pills spell out the type contract: what the stage above
        outputs, what the new stage must accept and emit, and what the stage below requires (or,
        at the end, that the chain's terminal type will change).
  \item Narrow the list with the search box or one of the category tabs: \ui{all} and the six
        categories \ui{preprocess}, \ui{encode}, \ui{detect}, \ui{cluster}, \ui{model} and
        \ui{control}. Each block card carries its glyph, name, type signature and a cost chip.
  \item Blocks that fit are listed first under \ui{Fits here}. Blocks that do not fit stay
        visible under \ui{Doesn't fit at this point}, disabled, each with the reason (for example
        ``needs Signal · here: Encoding''). \ui{show incompatible} hides them.
  \item Click a fitting card. The right-hand panel shows its description, its defaults, an
        estimated cost on this span, its side inputs, and warnings: which row will go stale,
        and whether the span is over the block's local ceiling.
  \item Press \ui{+ Insert}, or \ui{\textrightarrow{} Insert and open settings} to go straight
        to the new block's page.
\end{enumerate}
Below the last row a suggestion card names what fits at the end of the chain, with
\ui{Browse compatible} and a one-click \ui{+ Insert \ldots} for the first block that fits.

\screenshot{pages/analyse/analyse.chain--fixup-d-feature-blocks-in-insert-modal}{The
\ui{Insert a stage} window at the end of a SpanSet chain: the three type pills at the top, ``2 of
41 blocks fit'', and every other block greyed with its reason.}

\subsection*{Valid and invalid junctions}
Every edit is checked by the bridge (about 150\,ms after you stop changing things). Where a block
cannot take what the row above hands it, the \ui{+ insert} pill above it is replaced by a red
pill, for example ``03 Threshold to spans needs Scores · 02 Baseline removal emits Signal''. The
pill offers \ui{+ Insert \ldots{} here} (the first block that would repair the junction) and
\ui{Show blocks that fit}. While any junction is red the estimate chip reads ``1 invalid
junction'', the footer reads \ui{Chain is invalid}, and \ui{Run chain} is disabled.

\screenshot{pages/analyse/analyse.chain--1e-invalid-junction}{A red junction, here on the
built-in demonstration chain, with the two repairs it offers.}

\begin{plainterms}[What the types are]
A block's type signature reads like ``Signal \textrightarrow{} Scores'': it takes a signal and
emits a score for every moment. There are seven types: Signal, Scores, SpanSet (a list of
intervals), Encoding (letters or images), WindowSet (fixed windows), Grouping (clusters) and
Model. The type of the last block decides what the chain \emph{is}: a chain ending in a SpanSet
is a detection chain; one ending in a Model is a training chain. \Cref{ch:core-ideas} explains
each type.
\end{plainterms}

\subsection*{Run the chain}
\begin{enumerate}
  \item Check the estimate chip. It shows the cost from the first row that needs computing, for
        example ``\(\le\) 40\,s core est.\ · 01 \textrightarrow{} 03'', or ``all in the step cache''
        when nothing needs computing.
  \item Press \ui{Run chain}. Rows land as they finish. A running row shows the time spent and,
        when the block reports it, its own fraction done; a block that reports no fraction shows
        a moving bar and says so. Rows waiting their turn read ``waits for \ldots''.
  \item \ui{Cancel} stops the run. A block that checks for a cancel stops inside its stage; any
        other stops when its stage ends. The stages already finished are kept.
  \item When the run ends, the footer gives the result (for a SpanSet chain, how many spans), the
        job and database run numbers, how many rows were written to detections, and the time.
\end{enumerate}
Each row's badge says where its result came from: \texttt{cached · 0 s} (restored from the step
cache), \texttt{computed · now cached}, \texttt{stale}, \texttt{waiting}, \texttt{failed},
\texttt{blocked} (an earlier row failed) or \texttt{new} (never run).

\screenshot{pages/zzz_surrogates/analyse.chain--fixup-t-run-with-null-reads-detected-versus-surrogate}{A
four-block run with the surrogate null switched on: every row computed and now cached, each drawn
on the shared axis.}

\begin{watchout}[Long runs go to the cluster, not the laptop]
Before a run starts, the bridge estimates how long it will take. If the estimate is over the
Analyse local limit (Settings \textrightarrow{} Compute \& HPC; 15 minutes by default), the run
is \emph{refused} and a SLURM job script is written instead. A red card says how long the chain
would take, which stage is the slow one, the longest span that would run locally, and where the
script is; \ui{Copy script} copies it. If the cluster's own wall-time cap would still kill the
job, the card says that too. Separately, a block with a hard sample ceiling turns its row into an
amber card (``this run would fail locally''); there the \ui{Create SLURM script} and
\ui{Upload computed artifact} buttons are present but not yet wired. While a block's cost model is
still being timed on a new machine (``calibrating''), \ui{Run chain} waits.
\end{watchout}

\subsection*{A failed block}
When a block raises an error, its row is replaced by a red card: ``04 Drop detection failed after
\ldots'', the error message, the adapter and the recipe. The card has \ui{View log} (the full
traceback), \ui{Open settings} and \ui{Retry 04}. Rows below it read \texttt{blocked}. The footer
says ``No result'' and that nothing was written to detections. The rows above the failure stay
cached, so the toolbar button reads \ui{Retry from 04} and recomputes only from there.

\screenshot{pages/analyse/analyse.chain--1f-failed-block}{A failed block on the demonstration
chain: the error card with \ui{View log}, \ui{Open settings} and the retry button.}

\subsection*{Change something and re-run from a step}
Editing a block (on its own page), inserting or deleting one marks that row and every row below
it stale: the old result stays drawn under a ``last run shown · stale'' pill so you can compare.
The run button becomes \ui{Re-run from 0\(n\)}, and only the stale suffix is recomputed: the rows
above are read from the step cache.

\begin{plainterms}[The step cache]
Every row's result is stored under a fingerprint of everything that produced it: the span and
every setting of that block and of every block above it. Change row 03 and rows 01 and 02 still
have the same fingerprint, so they come back from the cache in no time. This is why tuning the
last block of a long chain is fast.
\end{plainterms}

To delete a block, use the bin icon on its row. A toast offers \ui{Undo}, and
\ui{Ctrl Z} restores the last deleted block. The grip, bypass and duplicate
icons on each row are present but not active on this page.

\subsection*{Run history}
\ui{History} lists the runs in the database for the current recording, plus the jobs live in this
bridge: name or number, the blocks, the span, the duration and how many detections each wrote.
\ui{Apply to source} loads a run's recipe into the chain (it replaces the unsaved chain); it is
disabled, with the reason, for a run on a different recording or with no stored recipe.
Discovery fan-outs and seed searches are counted at the foot rather than listed; open them in Jobs.

\subsection*{Import and save a template}
\ui{Import} opens the templates list, grouped by kind (Detection, Encoding, Training,
Interrogation) with a search box. Each row is marked \texttt{builtin} or \texttt{saved} and lists
its blocks. Picking one replaces the chain and forgets the last run's rows. Under it, a second list
of \emph{concept-frame chains} (marked \texttt{demo}) opens built-in demonstration chains with
simulated runs and example results; they are for learning the page, not for analysis.

\ui{Save template} asks for a name and saves the chain's blocks and settings as a template, which
is offered by \ui{Import} from then on; a detection template is also offered by Discovery's
\ui{Apply template}. The footer chip says what it
will be saved as (for example ``terminal SpanSet \textrightarrow{} detection template'').

\subsection*{The surrogate null}
The \ui{surrogate null} switch is off by default. Switched on, the next run is followed by its
paired \term{null}{null}: the whole chain is run again on several \term{surrogate}{surrogates} of
the span (the method and number of draws set in Settings \textrightarrow{} Nulls), and the footer
reads ``detected \(n\) · null expects \(m\) over \(k\) draws · \(x\times\) null''. The draws are a
count, never detections.

\pagefield{Where it leads}
\begin{itemize}
  \item A row's settings icon opens \apppage{analyse.block}.
  \item \ui{Pass \(n\) to Review} (footer, for a completed SpanSet chain) makes a Review queue over
        the detections the run wrote and opens it in \apppage{review.inspector}.
  \item \ui{Save window set} on a WindowSet row saves its windows for Library and Models.
  \item On a window-pool chain ending in Shape clustering, \ui{Train model} saves the template and
        opens \apppage{models.launch} with the pool filled in.
  \item \ui{Export run} writes a JSON report of the run. \ui{Analyse events} is not wired yet.
  \item A saved template is applied across many channels in \apppage{discovery.runs}.
\end{itemize}

\begin{underthehood}
Client: \code{webui/client/src/analyse/ChainPage.tsx} (\code{ChainPage}), with
\code{InsertStageModal.tsx}, \code{ChainRow.tsx}, \code{rowState.ts} (\code{deriveRows}: the row
badges), \code{useValidation.ts}, \code{HistoryPopover.tsx}, \code{TemplatesPopover.tsx},
\code{store.ts} (the run, streamed over server-sent events). Bridge: \code{webui/server/app.py}
routes \code{POST /api/chain/validate} (junctions, estimate, cache prediction, route),
\code{POST /api/chain/compatible} (what fits at a position, used by the insert window and the
junction repairs), \code{POST /api/runs} (refuses an invalid chain, an over-ceiling stage, or a
chain over the local limit, writing its SLURM job with \code{Working/hpc/job_export.py}
\code{export_job}), \code{GET /api/runs/{id}/log}; templates in
\code{webui/server/analyse_routes.py}. Routing: \code{webui/server/chain.py}
\code{route_local} and \code{analyse_ceiling_s}, falling back to
\code{Working.config.CLUSTER_ROUTING_CEILING_S} (900\,s).
\end{underthehood}

% ---------------------------------------------------------------------------------------------
\pagesection{analyse.block}{Block}{Settings and output together}{live}
The block page opens one block of the chain. The chain row shows only a block's result; this page
shows its process: what the block looked at, what it computed, and what it emitted, beside its
settings. You reach it from the settings icon on a row, or by clicking a block in the ribbon at the
top of the page.

\pagefield{When you'd come here}
\begin{itemize}
  \item To tune a block's settings and see the effect.
  \item To read a block's \emph{readout}: what it works out before a run (a noise level, a cut,
        a count of segments) and what it warns about.
  \item To look closely at a result: the detected spans one by one, a cluster tree, a model's
        accuracy per class.
\end{itemize}

\screenshot{pages/zz_analyse_run/analyse.block--fixup-h-spanset-cut-and-slideshow}{A threshold
block's page after a run: its cut drawn on the scores above, the spans it produced, and a
slideshow of them.}

\pagefield{Doing things here}
\begin{enumerate}
  \item The ribbon across the top shows the whole chain as chips, each with its status badge.
        Click a chip to move to that block; \ui{+ stage} goes back to the chain page to insert at
        the end; \ui{full chain} returns to the chain page.
  \item The large card on the left is the block's output. What it draws depends on the block's
        output type and input type, not on the block's name: a filter draws the signal before
        (grey) and after (blue); a threshold draws its cut on the scores and the spans it made;
        an encoding draws its letters or images; a clustering draws its tree.
  \item The \ui{Parameters} card on the right holds one control per setting: a segmented choice
        for up to four options, a menu for more, a slider and number box for a bounded number, a
        tick box for on/off. Under each control, ``= default'' or ``default \ldots'' says whether
        you have moved it; the \ui{i} beside a name shows its description. Some blocks (the
        bandpass filter) add preset buttons from the named bands in Settings.
  \item Change a setting. The block and every block after it go stale, and the footer reads
        \ui{Unapplied changes} with the cost of re-running. Press \ui{Re-run from 0\(n\)} (in the
        toolbar or the footer) to recompute from the first stale block; the page stays put and
        refreshes when the run ends.
  \item A block that declares a readout shows a \ui{Readout} card above its parameters. It is
        recomputed a quarter-second after each edit, on what a run would hand this block: the
        blocks before it are read from the step cache or re-run in memory, and nothing is written.
        Warnings are amber, errors red.
  \item \ui{Revert to defaults} puts every setting back to the block's default.
  \item \ui{Delete block} removes this block and returns to the chain page, where
        \ui{Ctrl Z} undoes it.
\end{enumerate}

\screenshot{pages/analyse/analyse.block--live}{The block page's layout: the ribbon, the output card,
the \ui{Parameters} card and the null card. With no source, the output card says it needs a span.}

\begin{watchout}[There is no per-parameter null sweep yet]
The concept design put a ``this parameter against the null'' sweep beside the settings. On the live
page that card only says that a null is drawn per chain run, not per block: switch on the
surrogate toggle and read the footer. Only the demonstration chains (opened from \ui{Import}
\textrightarrow{} concept-frame chains) draw a sweep, and theirs is example data.
\end{watchout}

\pagefield{Where it leads}
Back to \apppage{analyse.chain} to run the whole chain. A SpanSet block's slideshow can send its
spans to Review (the same queue as the chain page's \ui{Pass \(n\) to Review}); a WindowSet block
offers \ui{Save window set} (Library \textrightarrow{} Window sets). Each block's card in
\cref{ch:blocks} explains its settings.

\begin{underthehood}
Client: \code{webui/client/src/analyse/BlockPage.tsx} (\code{BlockPage}), \code{ParamsPanel.tsx}
(controls generated from the adapter's declared parameters), \code{Readout.tsx} and
\code{useReadout.ts}, and \code{views/registry.tsx}, which picks the drawing by output type and
conversion (\code{VIEWS}, \code{MODIFIERS}), mirrored on the server by
\code{webui/server/views.py}. Readout: \code{POST /api/blocks/derive}
(\code{webui/server/readout.py}). Parameter validation: \code{POST /api/chain/params}.
\end{underthehood}

% ---------------------------------------------------------------------------------------------
\pagesection{analyse.glyphs}{Algorithm glyphs}{The glyph registry}{preview}
Every block carries a small fixed picture, its \emph{glyph}, drawn once for the algorithm and
never from data. The same glyph appears at thumbnail size on the cards of the \ui{Insert a stage}
window and larger in its detail panel. This page lays all the glyphs out on one sheet with their
colour key: grey for the input or context, blue for what the block emits, green for what it found
or kept, amber for a cut or threshold, red for what it excluded, purple for an exemplar or second
input.

The page is a preview. It lists the blocks of the built-in demonstration chains (from
\code{webui/client/src/fixtures/analyse.ts}), not the live registry, and its header carries the
demo marker. The glyphs on the live \ui{Insert a stage} cards come from the same drawing code
(\code{webui/client/src/analyse/glyphs.tsx}), where a block with no glyph of its own falls back to
one drawn from its type signature.

\pagefield{When you'd come here}
Rarely: to see what a glyph means, or, when adding a block, to see which blocks still lack a glyph
of their own.

\screenshot{pages/analyse/analyse.glyphs--default}{The glyph sheet, grouped by kind of block, with
the colour key above it.}

\pagefield{Doing things here}
\begin{enumerate}
  \item Type in \ui{filter glyphs} to narrow the sheet by name or signature.
  \item Click a glyph to open its drawer: the glyph at three sizes, the block's description,
        registry name, category, whether it has a surrogate null, and which demonstration
        templates use it. \ui{Copy registry name} copies the name; \ui{Find it in Insert a stage}
        opens the insert window on a demonstration chain.
\end{enumerate}

\pagefield{Where it leads}
Back to \apppage{analyse.chain}. To add a glyph for a new block, see \cref{ch:extending}.

% ---------------------------------------------------------------------------------------------
\pagesection{analyse.interrogation}{Interrogation}{Source block: the family being analysed}{preview}
An interrogation chain starts not from a span of signal but from a \term{family}{family}: a group
of events the Library holds as similar shapes. Its source block decides which members of the
family the chain runs over; the blocks after it measure every member (01) and then plot the
measures against each other (02). This page is that source block.

\textbf{What is real here.} The families and their members are live: the bridge reads them from
the seed store of drop motifs (\code{DATA/library_seed/}), 16 spans holding 410 events, with the
detrended millivolt samples of every member. The \ui{prior run} tab of the source picker lists
real runs from the database. Everything that would \emph{run} or \emph{write} is simulated:
\ui{Run chain} plays a timed animation, \ui{Save template} keeps the template in browser memory
only (the toast says ``in memory (demo)''), and the clustering menu, the Review and Explore tabs of
the source picker and its \ui{prior run} rows are not wired.

\pagefield{When you'd come here}
To look at every member of one family at once, choose which of them belong in an analysis, and go
on to their slope or shape measures. It is the way into \apppage{analyse.interrogation.slope} and
\apppage{analyse.interrogation.aggregate}.

\screenshot{pages/interrogation/analyse.interrogation--default}{A family as the source block: one
slideshow card per member drawn from its stored samples, the scope matrix by recording and
channel, and the source settings.}

\pagefield{Doing things here}
\begin{enumerate}
  \item Open the source chip in the toolbar to pick a family (the \ui{family} tab). The header
        names the family and how many of its members are in scope.
  \item Narrow the members with the filters: \ui{adjudicated only}, \ui{exclude artifacts},
        \ui{distance \(\le\)}, \ui{recording} and \ui{channel}. The held-out recording is offered
        but disabled.
  \item Tick or untick a member on its card to put it in or out of scope; \ui{select all},
        \ui{none} and \ui{invert} act on all visible members. Excluding a member scopes this run
        only: the Library entry is not changed. The scope matrix on the right counts the members
        in scope per recording and channel.
  \item \ui{Members overlaid} draws a sample of ten members in scope, on a y axis measured from
        the traces drawn (never normalised). \ui{resample} draws another ten; \ui{align} lines
        them up on the onset or the trough; \ui{frame} switches between one context window for
        all and a \emph{sequence frame} that reaches back to the previous event's trough.
  \item \ui{Source settings} \textrightarrow{} \ui{context padding} sets the window drawn around
        every event on all three Interrogation pages; the default is the whole stored window.
  \item \ui{+ stage} on the chain card chooses what block 01 is: \ui{01 Resolve spans} (the
        store's slope analysis) or \ui{01 Event shape} (\blockref{interrogation.event_shape}).
  \item Click a member's card to open it in block 01.
\end{enumerate}

\pagefield{Where it leads}
\apppage{analyse.interrogation.slope} (block 01), \apppage{analyse.interrogation.aggregate}
(block 02), \ui{Open in Library} for the family's own page, and \ui{Sequences · slope rose} for the
steepest-slope rose across the events of one stored sequence.

\begin{underthehood}
Client: \code{webui/client/src/interrogation/SourcePage.tsx}, \code{SourcePicker.tsx},
\code{chrome.tsx}; reads in \code{webui/client/src/api/interrogation.ts} (\code{getSourceBlock},
\code{getSourceChoices}) mix live calls with constants from
\code{webui/client/src/fixtures/interrogation.ts} (the clustering label, source-setting options,
estimate and chain card). Bridge: \code{GET /api/interrogation/families} and
\code{GET /api/interrogation/families/{key}/members} in
\code{webui/server/interrogation_routes.py}. The run is \code{useSim}, a client-side timer.
\end{underthehood}

% ---------------------------------------------------------------------------------------------
\pagesection{analyse.interrogation.slope}{Interrogation: slope}{Slope analysis}{preview}
Block 01 of an interrogation chain: every member of the family turned into numbers. With
\ui{01 Resolve spans} the numbers are the seed store's slope measures; with \ui{01 Event shape}
they are the measures of \blockref{interrogation.event_shape} (amplitude, width, duration, full
width at half maximum, recovery time, rise time, steepest, onset and chord slope, peakedness).

\textbf{What is real here.} The events, their measures, the marks drawn on each event and the
rose of steepest slopes all come from the bridge, measured by the core; the browser computes none
of them. The rule selectors are disabled, because nothing on this page re-runs the measurement
yet; \ui{Run chain} and \ui{Apply} are simulated; \ui{CSV} is not wired.

\pagefield{When you'd come here}
To see how each number is defined on one event, check the measures across the family, and find
events the measurement could not handle.

\screenshot{pages/interrogation/analyse.interrogation.slope--default}{One event with the store's
marks (onset, steepest sample, trough), the strip of neighbouring events, and the per-event
table.}

\pagefield{Doing things here}
\begin{enumerate}
  \item The anatomy card draws one event with its marks: the detector's onset and trough and the
        steepest sample between them, with the tangent there. Step through events with the strip
        beneath it, or click a row of the table.
  \item \ui{Events overlaid} draws a seeded sample of the family with the current event in blue.
  \item \ui{Each fall as one angle} turns every fall into one angle (the arctangent of its
        steepest slope over a stated reference) and draws them as a rose with its circular
        statistics.
  \item The \ui{Rules} card lists the rules every number on the page was measured with: the onset
        rule, the trough rule, the steepest window and the slope noise estimate.
  \item The \ui{Per-event output} table has one row per member. A measure the core could not
        take is shown as a dash or ``not recovered'' (for an event that never came half-way back),
        never as zero.
\end{enumerate}

\begin{watchout}[Not measured is not zero]
An event that did not recover within the rule's bound has no recovery time. The page says ``not
recovered'' and leaves it out of every distribution, counting it beside the plot instead. Read a
recovery-time histogram with that count in mind.
\end{watchout}

\pagefield{Where it leads}
\apppage{analyse.interrogation.aggregate} for distributions and scaling; back to
\apppage{analyse.interrogation} to change the scope.

\begin{underthehood}
\code{webui/client/src/interrogation/SlopePage.tsx}; \code{getSlopeBlock} in
\code{webui/client/src/api/interrogation.ts}; bridge routes
\code{GET /api/interrogation/families/{key}/slope} and \code{.../shape}, measuring with
\code{Working/Detection/drop_motifs/gradients.py} and the \code{interrogation.event_shape} block.
\end{underthehood}

% ---------------------------------------------------------------------------------------------
\pagesection{analyse.interrogation.aggregate}{Interrogation: aggregate}{Aggregate}{preview}
Block 02 of an interrogation chain. It takes whatever measures block 01 emits and draws them:
a histogram per measure, a scaling plot of one measure against another on log axes with a fitted
exponent, and a timeline of when each event happened in each recording. The block is generic: its
plots are wired to the upstream block's declared features, so it serves the slope analysis and the
event-shape block alike.

\textbf{What is real here.} The measures are the core's (as on block 01); the fit is computed on
the points drawn. The page draws \emph{no null}: none has been built for these plots, and the page
says so on its face. \ui{CSV}, \ui{Figures} and \ui{add scaling pair} are not wired; sending an
outlier to Review is simulated.

\pagefield{When you'd come here}
To see how a family's events are distributed, and whether one measure scales with another (for
example whether deeper drops are steeper).

\screenshot{pages/interrogation/analyse.interrogation.aggregate--default}{Histograms of the
family's measures, the scaling plot with its fitted exponent, and the parameters card.}

\pagefield{Doing things here}
\begin{enumerate}
  \item Pick the pair of measures for the scaling plot, and switch its axes between log and
        linear. The fitted exponent \(\beta\) is shown with a 95\,\% interval; with fewer than 12
        points the page declines to fit.
  \item Colour the points by recording to see whether one recording drives the fit.
  \item \ui{Feature wiring} shows which upstream feature feeds each plot slot, and lets you
        rewire or reset it.
  \item The \ui{Parameters} card sets the binning rule, the interval definition, how outliers are
        handled and the purity check. These change only the views: nothing here is stored, and
        nothing upstream goes stale.
  \item \ui{How each measure is defined} prints the rules the core measured with.
\end{enumerate}

\begin{watchout}[A fitted exponent is a description, not a test]
With no null drawn, \(\beta\) describes these events; it does not show that a relationship is
there rather than chance. Use Discovery's nulls for claims (\cref{ch:discovery}).
\end{watchout}

\pagefield{Where it leads}
Back to \apppage{analyse.interrogation.slope} for single events.

\begin{underthehood}
\code{webui/client/src/interrogation/AggregatePage.tsx}, \code{features.ts} (\code{featureOf}),
\code{EventTimeline.tsx} (heights from \code{POST /api/interrogation/trend});
\code{GET /api/interrogation/families/{key}/aggregate}. The fit is \code{fitLogLog} in
\code{webui/client/src/fixtures/interrogation.ts}: ordinary least squares on log10 of both
measures, computed in the browser.
\end{underthehood}

% ---------------------------------------------------------------------------------------------
\pagesection{analyse.training}{Training}{Built on one channel, applied across the corpus}{preview}
The Training pages are a design preview of a \emph{training chain}: windows cut from one channel,
features measured per window, windows clustered, encoded as images, and a model trained on them.
Every number on these six pages is example data from \code{webui/client/src/fixtures/training.ts};
nothing on them calls the bridge, and every save, queue or job they make lives only in browser
memory (the toasts say ``in memory (demo)'').

\begin{watchout}[Use the live chain instead]
The same work is done for real on \apppage{analyse.chain}: import the \texttt{windows\_model}
template and run it on a span, or start from a Window pool (\ui{Start from a Window pool}) to
cluster saved window sets with \blockref{catalogue.shape_cluster}, then press \ui{Train model} to
go to \apppage{models.launch}. The preview's cluster page itself points there.
\end{watchout}

\pagefield{When you'd come here}
To see the intended design of a training chain: which checks guard against a model learning from
its own test data, and how a template built on one channel is handed to Models.

\screenshot{pages/training/analyse.training--default}{The preview training chain: five rows on one
channel, ending in a Model, with the hand-off to Models in the footer.}

\pagefield{Doing things here}
\begin{enumerate}
  \item Read the five rows: \ui{01 Sliding windows}, \ui{02 Window matrix}, \ui{03 Cluster},
        \ui{04 Encode}, \ui{05 Model}. The chain is a training chain because its last block emits
        a Model.
  \item Switch the source to \emph{human-labelled windows} to see the illustrative variant, in
        which the sliding-windows stage is absent and the stages renumber.
  \item Open a row to see its block page (below).
\end{enumerate}

\pagefield{Where it leads}
The five block pages below; \ui{Train in Models} opens \apppage{models.launch}.

\begin{underthehood}
\code{webui/client/src/training/} (\code{ChainPage.tsx} and one page per block); reads in
\code{webui/client/src/api/training.ts}, every one \code{demo(...)} over a fixture.
\end{underthehood}

\pagesection{analyse.training.windows}{Training: sliding windows}{Windows, stride and the blocked split}{preview}
Block 01 of the preview: the channel cut into fixed windows, and the windows split into training,
validation and test blocks by time. Its point is the guard against leakage: the gap between blocks
is a setting, checked to be at least one window long, and a random split can be chosen but is
marked as leaking. The live block is \blockref{preprocessing.sliding_windows}; the live split is
set on \apppage{models.launch}.

\begin{plainterms}[Why the split is by time]
Neighbouring windows overlap and share samples. Split them at random and a model is tested on
windows it has half seen in training, like a student marked on the questions they revised. Cutting
the recording into blocks of time, with a gap between them, keeps the test honest. This is a
\term{blocked-split}{blocked split}.
\end{plainterms}

\screenshot{pages/training/analyse.training.windows--default}{Example windows on the channel, the
train, validation and test blocks with their gaps, and the checks card.}

\pagefield{When you'd come here}
To see how the blocked split is meant to be set and checked.

\pagefield{Doing things here}
Change the window length, stride, gap and split ratio in \ui{Parameters}; the \ui{Checks} card says
which conditions hold. All of it is example data.

\pagefield{Where it leads}
\apppage{analyse.training.matrix}.

\pagesection{analyse.training.matrix}{Training: window matrix}{Features per window}{preview}
Block 02 of the preview: a table of features measured for every window, in feature groups that
can be ticked on and off. Two groups whose scores come from models trained on the manual verdicts
are off by default, and ticking one raises a leakage warning, because the clustering is later
compared against those same verdicts. The page also shows how a heavy matrix would go to the
cluster (\ui{Create SLURM script}) and come back (\ui{Upload matrix}); both are simulated. The live
block is \blockref{preprocessing.window_matrix}.

\screenshot{pages/training/analyse.training.matrix--default}{The example feature matrix, one row
per window, under the signal, with the feature groups on the right.}

\pagefield{When you'd come here}
To see which feature groups a window matrix carries and which of them would leak labels.

\pagefield{Doing things here}
Tick feature groups; open one to see its columns; read a window's values in the readout and open
it in \ui{04 Encode}.

\pagefield{Where it leads}
\apppage{analyse.training.cluster}.

\pagesection{analyse.training.cluster}{Training: cluster}{Clusters and choosing k}{preview}
Block 03 of the preview: the windows grouped by feature similarity into a tree, cut into \(k\)
classes. Its second view, \ui{choose k}, sets out the evidence for the number of classes and lets
you lock the selection criterion \emph{before} looking, because a criterion chosen afterwards makes
\(k\) the finding rather than the method. The page's own callout points to the live dendrogram:
the Shape clustering block's page (\blockref{catalogue.shape_cluster}), reached from a Window pool
chain on \apppage{analyse.chain}.

\screenshot{pages/training/analyse.training.cluster--choose-k}{The choose-\(k\) view: the
criterion cards, the sweep over \(k\), and the lock.}

\pagefield{When you'd come here}
To see the intended guard on choosing \(k\); for real clustering, use the live chain.

\pagefield{Doing things here}
Drag the cut on the example tree, merge classes, open \ui{choose k}, lock or unlock the criterion.
None of it is stored.

\pagefield{Where it leads}
\apppage{analyse.training.encode}.

\pagesection{analyse.training.encode}{Training: encode}{Image encodings}{preview}
Block 04 of the preview: each window turned into the images a classifier would see (Gramian and
recurrence encodings), one window at a time, by class. It also states the rule that an encoder set
or image size not already registered becomes a new encoder version, so models trained earlier keep
theirs. The live encodings are blocks such as \blockref{catalogue.gramian_gasf},
\blockref{catalogue.gramian_gadf} and \blockref{catalogue.gramian_recurrence}.

\screenshot{pages/training/analyse.training.encode--default}{One example window and its image
encodings, with the encoder set and its cost tiles.}

\pagefield{When you'd come here}
To see what the classifier is meant to be given.

\pagefield{Doing things here}
Page through windows, browse by class, tick encodings and change the image size.

\pagefield{Where it leads}
\apppage{analyse.training.model}.

\pagesection{analyse.training.model}{Training: model}{The model stage}{preview}
Block 05 of the preview, the terminal block. It does not train: it says what a trial on this one
channel would cost, which labels the model would learn from (cluster classes, manual verdicts or
both, paired), what stands in the way (the held-out recording is locked), and hands the template to
Models, where a template is applied across channels. The live classifier is
\blockref{catalogue.classifier}.

\screenshot{pages/training/analyse.training.model--default}{The model card with its label source,
the checks before training and the hand-off to Models.}

\pagefield{When you'd come here}
To see how a training template is meant to be handed to Models.

\pagefield{Doing things here}
Choose where the labels come from; read the checks; press \ui{Train in Models}.

\pagefield{Where it leads}
\apppage{models.launch}.
